
\chapter{Especificações completas dos modelos}
\label{ap:especificacoes}

Este apêndice reúne, para cada desenho inferencial de que dependem as
afirmações do Capítulo~\ref{cap:resultados}, a equação estimada, a definição e
a escala das variáveis, a estrutura de efeitos fixos, o agrupamento dos
erros-padrão, o número de observações e de unidades, e o resultado de \emph{todas}
as especificações rodadas em cada grade --- e não apenas das que o texto
comenta. O objetivo é que os modelos possam ser reconstruídos a partir deste
documento, sem recurso ao repositório. São sete blocos: a precedência temporal
entre as séries regionais; o teste espacial do deslocamento; a dependência
espacial nas formas SAR e SEM; o desenho \emph{shift-share} do \emph{drive}
comum, com a sua bateria de placebos e a corrida entre exposições rivais; a
grade confirmatória completa; o teto de oferta, com a grade inteira de
tratamentos da variável de depleção; e o desenvolvimento municipal.

\textbf{O que não está aqui, e por quê.} Três coisas. As \textbf{estatísticas
das quebras estruturais} já têm quadro próprio no corpo do texto
(Quadro~\ref{quadro:quebras}), e repeti-las aqui seria duplicação. O
\textbf{ajuste de mistura da idade da pastagem} não é regressão: seus
parâmetros são médias, dispersões e pesos das duas componentes, e estão
enunciados na própria Seção~\ref{sec:res-perna2}, com a figura que os sustenta.
E a \textbf{grade exploratória} do \emph{drive} comum --- 192 combinações,
contra as 14 da confirmatória --- fica de fora por decisão de desenho, não de
espaço: ela existe para gerar hipóteses sob controle de multiplicidade, e
nenhuma afirmação do texto se apoia nela. Quem quiser conferi-la encontra o
arquivo nomeado na Seção~\ref{sec:met-repro}.

As tabelas deste apêndice são geradas por \texttt{apendice/gerar\_apendice.py} a
partir dos mesmos arquivos que os pipelines gravam; nenhum coeficiente foi
transcrito à mão.

\section{Convenções comuns}
\label{ap:convencoes}

\textbf{Unidade e painel.} As regressões deste apêndice correm sobre o painel de 166
Áreas Mínimas Comparáveis de Goiás, 1985--2024 (\(166 \times 40 = 6.640\)
células), descrito na Seção~\ref{sec:met-amc}. Quando uma fonte não cobre a
janela inteira, o \(n\) reportado na tabela é o efetivo após exclusão de
ausentes por lista.

\textbf{Primeira diferença (D7).} O desfecho é sempre a variação anual, e não o
nível: \(\Delta y_{it} = y_{it} - y_{i,t-1}\). É essa transformação que desarma
o vazamento entre regressor e regressando documentado na
Seção~\ref{sec:met-painel}.

\textbf{Padronização.} As exposições entram em \emph{escore-z} calculado sobre
a seção transversal das 166 unidades \(\tilde{E}_i = (E_i - \bar{E})/\sigma_E\),
com \(\sigma\) populacional. Os coeficientes de interação lêem-se, portanto,
como efeito de um desvio-padrão de exposição.

\textbf{Efeitos fixos e erros-padrão.} Salvo indicação em contrário, todas as
especificações trazem efeito fixo de unidade \emph{e} de ano (2FE) e
erros-padrão agrupados nas duas dimensões. O agrupamento efetivamente aplicado
é reportado em cada tabela, porque em uma especificação ele recai para
agrupamento só por unidade quando a matriz de covariância bidimensional não é
positiva-definida.

\textbf{Duas réguas de inferência.} Onde o desenho é do tipo
\emph{shift-share}, o \(p\) agrupado e o \(p\) de permutação circular do
\emph{shifter} são reportados lado a lado. Eles não são comparáveis entre si: o
primeiro trata cada ano como realização independente do choque, e o segundo
não. A régua defensável é a segunda, pela razão exposta na
Seção~\ref{sec:met-instrumental}.


\section{Precedência temporal entre as séries regionais}
\label{ap:precedencia}

Antes do teste espacial, a hipótese do empurrão foi examinada no tempo: se a
lavoura do Sul puxa a pastagem do Norte, o passado de uma deve ajudar a prever
a outra. Este é o único bloco do apêndice que \emph{não} usa o painel de AMCs.
Ele corre sobre \textbf{séries anuais agregadas por região}, o que dá cerca de
38 observações --- e é dessa escassez, e não do desenho, que vem a maior parte
da limitação de poder.

A leitura do teste de Granger tem de ser literal: ele mede \emph{precedência
preditiva}, não causalidade. E, aplicado a séries integradas, fabrica
precedência inexistente --- foi o que a decisão D16 registrou. Por isso a
Tabela~\ref{tab:precedencia} reporta, lado a lado, o \(p\) clássico e o \(p\)
de um teste de Wald com erros-padrão robustos a heterocedasticidade e
autocorrelação, na receita de Toda-Yamamoto, que permanece válido sob
integração.

A direção \emph{reversa} está na tabela de propósito. Ela é o achado que a
decisão D16 classificou como \textbf{artefato espúrio}: o sinal forte que
aparece de norte para sul não sobrevive ao tratamento da integração como
relação de longo prazo, e é reportado aqui para que o leitor veja o que o
diagnóstico teve de descartar.


\begin{table}[htb]
\centering
\caption[Precedência temporal entre séries regionais]{Precedência temporal entre as séries regionais, por defasagem.}
\label{tab:precedencia}
\small
\begin{tabular}{lcccccr}
\toprule
& \textbf{Def.} & \textbf{$p$ Granger} & \textbf{$p$ Wald HAC} & \textbf{$\hat{\beta}_1$} & \textbf{$p(\hat{\beta}_1)$} & \textbf{$n$} \\
\midrule
\multicolumn{7}{l}{\textit{Sul$\rightarrow$Norte: $\Delta$Agric\_Sul $\rightarrow$ $\Delta$Pasto\_Norte}} \\
 & 1 & $0{,}971$ & $0{,}959$ & $+0{,}0022$ & $0{,}959$ & 38 \\
 & 2 & $0{,}760$ & $0{,}703$ & $-0{,}0257$ & $0{,}669$ & 37 \\
 & 3 & $0{,}822$ & $0{,}518$ & $-0{,}0202$ & $0{,}744$ & 36 \\
\multicolumn{7}{l}{\textit{REVERSO:   $\Delta$Pasto\_Norte $\rightarrow$ $\Delta$Agric\_Sul}} \\
 & 1 & $<0{,}001$ & $<0{,}001$ & $+0{,}4394$ & $<0{,}001$ & 38 \\
 & 2 & $0{,}075$ & $0{,}005$ & $+0{,}3494$ & $0{,}339$ & 37 \\
 & 3 & $0{,}179$ & $0{,}006$ & $+0{,}2306$ & $0{,}554$ & 36 \\
\multicolumn{7}{l}{\textit{REVERSO bov: $\Delta$Bovinos\_Norte $\rightarrow$ $\Delta$Agric\_Sul}} \\
 & 1 & $0{,}723$ & $0{,}650$ & $-0{,}0151$ & $0{,}647$ & 38 \\
 & 2 & $0{,}192$ & $0{,}407$ & $-0{,}0323$ & $0{,}392$ & 37 \\
 & 3 & $0{,}299$ & $0{,}346$ & $-0{,}0425$ & $0{,}308$ & 36 \\
\multicolumn{7}{l}{\textit{Sul$\rightarrow$Norte bov: $\Delta$Agric\_Sul $\rightarrow$ $\Delta$Bovinos\_Norte}} \\
 & 1 & $0{,}325$ & $0{,}180$ & $+0{,}6209$ & $0{,}171$ & 38 \\
 & 2 & $0{,}593$ & $0{,}242$ & $+0{,}6756$ & $0{,}095$ & 37 \\
 & 3 & $0{,}723$ & $0{,}332$ & $+0{,}9161$ & $0{,}095$ & 36 \\
\bottomrule
\end{tabular}
\fontefig{elaboração própria.}
\notafig{séries anuais agregadas por região, em primeira diferença; ``Def.'' é o número de defasagens do modelo. $p$ Granger é o teste-F clássico e $p$ Wald HAC o mesmo teste com covariância robusta a heterocedasticidade e autocorrelação; $\hat{\beta}_1$ é o coeficiente da primeira defasagem da série candidata a precedente. O bloco \textsc{reverso} é reportado porque foi ele que originou a decisão D16, e não porque sustente afirmação.}
\end{table}

\section{Teste espacial do deslocamento local (Perna 3a)}
\label{ap:slx}

A hipótese é a de que a lavoura que avança numa unidade empurre a pastagem das
unidades \emph{ao sul} dela. O desenho é um modelo espacial de defasagem dos
\emph{regressores} (SLX), estimado em primeira diferença com efeitos fixos de
unidade e de ano:

\begin{equation}
\Delta y_{it} = \beta\,\Delta x_{it}
             + \theta\,(W\Delta x)_{it}
             + \alpha_i + \lambda_t + \varepsilon_{it},
\label{eq:slx}
\end{equation}

\noindent em que \(\Delta y_{it}\) é a variação anual da área de pastagem (ou
do efetivo bovino, nas linhas assim indicadas) da unidade \(i\) no ano \(t\);
\(\Delta x_{it}\) é a variação anual da área de lavoura \emph{na própria
unidade}; \(\alpha_i\) e \(\lambda_t\) são os efeitos fixos; e \((W\Delta
x)_{it}\) é a média da variação de lavoura nas unidades vizinhas selecionadas
por \(W\). O parâmetro de interesse é \(\theta\), não \(\beta\): \(\beta\)
mede substituição \emph{dentro} da unidade, que não é deslocamento.

\textbf{Matriz de pesos.} \(W\) é \(166 \times 166\), construída sobre os oito
vizinhos mais próximos de cada unidade por distância euclidiana entre
centroides em EPSG:5880, e \textbf{padronizada por linha}. A direção é imposta
por filtro: em \(W_{\mathrm{sul}}\) só permanecem, entre os oito, os vizinhos de
centroide mais ao sul; em \(W_{\mathrm{norte}}\), os mais ao norte. A diagonal é
nula. \(W_{\mathrm{norte}}\) é \emph{placebo}: se o mecanismo for o empurrão da
lavoura do sul, o termo de vizinhança deve aparecer com \(W_{\mathrm{sul}}\) e
não com \(W_{\mathrm{norte}}\).

\textbf{Grade de especificações.} A grade é o produto de três réguas de
lavoura --- a classe Agricultura do MapBiomas, a união dela com o Mosaico de
Usos (o \emph{bracket} da decisão D26) e a área plantada de soja do SIDRA, que
é imune ao classificador --- por duas janelas --- a série plena e a truncada em
2019, que exclui o trecho de deriva do Mosaico --- por três modelos, sendo um
deles o placebo com $W_{\mathrm{norte}}$. São dezoito especificações e trinta e
seis coeficientes, já que cada uma estima o termo local e o de vizinhança; a
Tabela~\ref{tab:slx} traz os dezoito termos de vizinhança, que são os de
interesse.


\begin{table}[htb]
\centering
\caption[Termo de vizinhança do teste espacial]{Termo de vizinhança do teste espacial do deslocamento, nas dezoito especificações.}
\label{tab:slx}
\small
\begin{tabular}{llccccr}
\toprule
\textbf{Régua de lavoura} & \textbf{Desfecho} & \textbf{$W$} & \textbf{$\hat{\theta}$} & \textbf{EP} & \textbf{$p$} & \textbf{$n$} \\
\midrule
\multicolumn{7}{l}{\textit{Janela plena (1985--2024)}} \\
Agricultura & $\Delta$ pastagem & $W_{\mathrm{sul}}$ & $-0{,}157$ & $0{,}068$ & $0{,}020$ & 6.474 \\
Agricultura & $\Delta$ pastagem & $W_{\mathrm{norte}}$ (placebo) & $+0{,}182$ & $0{,}142$ & $0{,}203$ & 6.474 \\
Agricultura & $\Delta$ bovinos & $W_{\mathrm{sul}}$ & $-0{,}124$ & $0{,}110$ & $0{,}261$ & 6.474 \\
Agric. $\cup$ Mosaico & $\Delta$ pastagem & $W_{\mathrm{sul}}$ & $-0{,}050$ & $0{,}082$ & $0{,}545$ & 6.474 \\
Agric. $\cup$ Mosaico & $\Delta$ pastagem & $W_{\mathrm{norte}}$ (placebo) & $-0{,}017$ & $0{,}048$ & $0{,}717$ & 6.474 \\
Agric. $\cup$ Mosaico & $\Delta$ bovinos & $W_{\mathrm{sul}}$ & $-0{,}018$ & $0{,}147$ & $0{,}902$ & 6.474 \\
Soja (SIDRA) & $\Delta$ pastagem & $W_{\mathrm{sul}}$ & $-0{,}012$ & $0{,}019$ & $0{,}526$ & 3.960 \\
Soja (SIDRA) & $\Delta$ pastagem & $W_{\mathrm{norte}}$ (placebo) & $-0{,}054$ & $0{,}025$ & $0{,}032$ & 3.960 \\
Soja (SIDRA) & $\Delta$ bovinos & $W_{\mathrm{sul}}$ & $-0{,}024$ & $0{,}075$ & $0{,}752$ & 3.960 \\
\multicolumn{7}{l}{\textit{Janela truncada (1985--2019)}} \\
Agricultura & $\Delta$ pastagem & $W_{\mathrm{sul}}$ & $-0{,}122$ & $0{,}070$ & $0{,}083$ & 5.644 \\
Agricultura & $\Delta$ pastagem & $W_{\mathrm{norte}}$ (placebo) & $+0{,}141$ & $0{,}114$ & $0{,}215$ & 5.644 \\
Agricultura & $\Delta$ bovinos & $W_{\mathrm{sul}}$ & $-0{,}096$ & $0{,}121$ & $0{,}429$ & 5.644 \\
Agric. $\cup$ Mosaico & $\Delta$ pastagem & $W_{\mathrm{sul}}$ & $-0{,}075$ & $0{,}094$ & $0{,}421$ & 5.644 \\
Agric. $\cup$ Mosaico & $\Delta$ pastagem & $W_{\mathrm{norte}}$ (placebo) & $-0{,}014$ & $0{,}057$ & $0{,}805$ & 5.644 \\
Agric. $\cup$ Mosaico & $\Delta$ bovinos & $W_{\mathrm{sul}}$ & $-0{,}012$ & $0{,}164$ & $0{,}941$ & 5.644 \\
Soja (SIDRA) & $\Delta$ pastagem & $W_{\mathrm{sul}}$ & $-0{,}025$ & $0{,}015$ & $0{,}094$ & 3.222 \\
Soja (SIDRA) & $\Delta$ pastagem & $W_{\mathrm{norte}}$ (placebo) & $-0{,}053$ & $0{,}027$ & $0{,}052$ & 3.222 \\
Soja (SIDRA) & $\Delta$ bovinos & $W_{\mathrm{sul}}$ & $-0{,}007$ & $0{,}088$ & $0{,}934$ & 3.222 \\
\bottomrule
\end{tabular}
\fontefig{elaboração própria.}
\notafig{estimação por mínimos quadrados com efeitos fixos de unidade e de ano; erros-padrão agrupados por unidade. O coeficiente reportado é o do termo de vizinhança $\hat{\theta}$ da Equação~\ref{eq:slx}; o termo local $\hat{\beta}$ é omitido porque mede substituição dentro da unidade, que não é deslocamento. O $n$ menor nas linhas da soja reflete o início mais tardio da série municipal do SIDRA.}
\end{table}

\section{Dependência espacial: modelos de defasagem e de erro}
\label{ap:sarsem}

As unidades vizinhas não são independentes, e ignorá-lo estreita a barra de
erro. Três modelos de painel foram, por isso, reestimados em três formas: por
mínimos quadrados com efeitos fixos; com defasagem espacial da \emph{variável
dependente} (SAR); e com dependência no \emph{termo de erro} (SEM). A
comparação responde a uma pergunta só --- se os coeficientes de interesse
sobrevivem quando a dependência é modelada --- e a resposta consta na última
coluna da Tabela~\ref{tab:sarsem}.

A matriz \(W\) aqui é de contiguidade \emph{queen} entre as 166 unidades,
padronizada por linha, e não a matriz direcional de oito vizinhos da
Seção~\ref{ap:slx}: aqui o objeto é a dependência genérica, e não a direção
sul-norte. Os testes do multiplicador de Lagrange, em suas versões simples e
robusta, indicam qual das duas formas o dado prefere.


\begin{table}[htb]
\centering
\caption[Dependência espacial nos três modelos de painel]{Coeficientes de interesse sob mínimos quadrados, SAR e SEM.}
\label{tab:sarsem}
\footnotesize
\begin{tabular}{lp{3.3cm}cccccll}
\toprule
& \textbf{Especificação} & \textbf{$\hat{\beta}_{MQ}$} & \textbf{$\hat{\beta}_{SAR}$} & \textbf{$\hat{\beta}_{SEM}$} & \textbf{$\hat{\rho}$} & \textbf{$\hat{\lambda}$} & \textbf{Forma} & \textbf{Sobrev.} \\
\midrule
M1 & Intensificação: $\Delta$agricultura $\sim$ $\Delta$ VA agro & $-0{,}0047$ & $-0{,}0043$ & $-0{,}0046$ & $0{,}353$ & $0{,}368$ & error & sim \\
M2 & Crédito$\rightarrow$pasto: $\Delta$pastagem $\sim$ $\Delta$ SICOR + $\Delta$ VA agro & $-0{,}0040$ & $-0{,}0035$ & $-0{,}0034$ & $0{,}398$ & $0{,}401$ & lag & sim \\
M3 & Substituição local: $\Delta$pastagem $\sim$ $\Delta$agricultura & $-0{,}5458$ & $-0{,}4774$ & $-0{,}5356$ & $0{,}391$ & $0{,}414$ & error & sim \\
M3 & Substituição local: $\Delta$pastagem $\sim$ $\Delta$agricultura & $-0{,}5458$ & $-0{,}4813$ & $-0{,}5491$ & $0{,}532$ & $0{,}558$ & error & sim \\
\bottomrule
\end{tabular}
\fontefig{elaboração própria.}
\notafig{$W$ de contiguidade \emph{queen} entre as 166 unidades, padronizada por linha. $\hat{\rho}$ é o parâmetro de defasagem espacial do SAR e $\hat{\lambda}$ o de dependência do erro no SEM; ``Forma'' é a preferida pelos testes do multiplicador de Lagrange em versão robusta. ``Sobrev.'' indica se o sinal e a significância do coeficiente de interesse resistem à modelagem da dependência. Todos os coeficientes vêm de especificações em primeira diferença com efeitos fixos.}
\end{table}

\section{\emph{Drive} comum: desenho \emph{shift-share} (Perna 3b)}
\label{ap:shiftshare}

A hipótese é a de que um choque macroeconômico nacional atinja de modo
diferente lugares com exposição diferente. A especificação estimada é de
interação pura, em primeira diferença e com dois conjuntos de efeitos fixos:

\begin{equation}
\Delta y_{it} = \gamma\,\bigl(\tilde{s}_{t-1} \times \tilde{E}_i\bigr)
             + \alpha_i + \lambda_t + \varepsilon_{it},
\label{eq:shiftshare}
\end{equation}

\noindent em que \(\Delta y_{it}\) é a variação anual do efetivo bovino da
unidade \(i\) (em milhões de cabeças); \(\tilde{s}_{t-1}\) é o \emph{shifter}
nacional --- a variação padronizada do índice de taxa de câmbio efetiva real,
defasada um ano ---; e \(\tilde{E}_i\) é a exposição local padronizada, fixa no
tempo.

Três propriedades da equação merecem registro, porque determinam o que ela pode
e não pode responder. Primeira: \textbf{não há termos principais}. O
\emph{shifter} é comum a todas as unidades num dado ano e é absorvido por
\(\lambda_t\); a exposição é fixa no tempo e é absorvida por \(\alpha_i\). O
que resta identificado é apenas o \emph{gradiente} do efeito --- se o choque
morde mais onde a exposição é maior --- e nunca o seu nível. Segunda: por
isso mesmo, \(\gamma\) não mede o efeito do câmbio sobre o rebanho, e nenhum
número deste apêndice autoriza essa leitura. Terceira: o número de realizações
independentes do choque é o número de anos, cerca de 38, e não o número de
células do painel; nenhuma quantidade adicional de unidades espaciais o
aumenta. É essa a razão de o \(p\) agrupado e o \(p\) de permutação divergirem.

\subsection{Inferência por permutação}

A permutação circular reembaralha o \emph{shifter} ao longo dos anos,
preservando a sua autocorrelação, e recalcula \(\hat{\gamma}\) em cada
reembaralhamento sobre o desfecho duplamente centrado. O \(p\) reportado é a
proporção de reembaralhamentos que produzem estatística ao menos tão extrema
quanto a observada. Com 38 anos, o menor \(p\) atingível é da ordem de
\(1/38 \approx 0{,}026\), o que impõe um piso à resolução da régua.

\subsection{A corrida entre exposições}

A bateria de placebos do Capítulo~\ref{cap:resultados} é toda de \emph{desfecho}:
ela pergunta se o sinal aparece onde não deveria. A
Tabela~\ref{tab:horserace} responde a outra pergunta, que é a decisiva para uma
\emph{share} espacial --- se a exposição escolhida é a certa --- pondo as
exposições rivais na mesma regressão. A latitude entra como confundidor puro:
ela mede ``quão ao norte'' sem nenhum conteúdo agronômico.


\begin{table}[htb]
\centering
\caption[Corrida entre exposições rivais]{Corrida entre exposições rivais no desenho \emph{shift-share}.}
\label{tab:horserace}
\footnotesize
\begin{tabular}{lp{3.5cm}p{3.1cm}ccccc}
\toprule
\textbf{Esp.} & \textbf{Especificação} & \textbf{Exposição do termo} & \textbf{$\hat{\gamma}$} & \textbf{EP} & \textbf{$p_{\mathrm{agr}}$} & \textbf{$p_{\mathrm{perm}}$} & \textbf{$R^2_w$} \\
\midrule
S1 & aptidão sozinha (reproduz o \#52) & Aptidão edafoclim. (Embrapa) & $-0{,}033$ & $0{,}015$ & $0{,}026$ & $0{,}132$ & $0{,}0012$ \\
S2 & latitude sozinha & Latitude (confundidor puro) & $+0{,}051$ & $0{,}021$ & $0{,}015$ & $0{,}026$ & $0{,}0029$ \\
S3 & acesso sozinho & Distância ao núcleo de lavoura 1985 & $+0{,}028$ & $0{,}018$ & $0{,}118$ & $0{,}184$ & $0{,}0009$ \\
S4 & aptidão + latitude & Aptidão edafoclim. (Embrapa) & $-0{,}012$ & $0{,}012$ & $0{,}302$ & $0{,}474$ & $0{,}0030$ \\
S4 & aptidão + latitude & Latitude (confundidor puro) & $+0{,}046$ & $0{,}021$ & $0{,}028$ & $0{,}053$ & $0{,}0030$ \\
S5 & aptidão + acesso & Aptidão edafoclim. (Embrapa) & $-0{,}024$ & $0{,}016$ & $0{,}139$ & $0{,}132$ & $0{,}0013$ \\
S5 & aptidão + acesso & Distância ao núcleo de lavoura 1985 & $+0{,}015$ & $0{,}021$ & $0{,}482$ & $0{,}395$ & $0{,}0013$ \\
S6 & aptidão + latitude + acesso & Aptidão edafoclim. (Embrapa) & $-0{,}026$ & $0{,}016$ & $0{,}111$ & $0{,}132$ & $0{,}0039$ \\
S6 & aptidão + latitude + acesso & Latitude (confundidor puro) & $+0{,}081$ & $0{,}042$ & $0{,}052$ & $0{,}079$ & $0{,}0039$ \\
S6 & aptidão + latitude + acesso & Distância ao núcleo de lavoura 1985 & $-0{,}051$ & $0{,}043$ & $0{,}236$ & $0{,}237$ & $0{,}0039$ \\
\bottomrule
\end{tabular}
\fontefig{elaboração própria.}
\notafig{desfecho: variação anual do efetivo bovino. Efeitos fixos de unidade e de ano em todas as linhas; erros-padrão agrupados por unidade e por ano ($n = 6.308$ observações, 166 unidades). $p_{\mathrm{agr}}$ é o $p$ do erro-padrão agrupado e $p_{\mathrm{perm}}$ o da permutação circular do \emph{shifter}; as duas réguas não são comparáveis entre si, e a segunda é a defensável. $R^2_w$ é o $R^2$ \emph{within}. Quando uma especificação traz mais de uma exposição, cada linha reporta o termo de interação da exposição nomeada, com as demais mantidas na mesma regressão.}
\end{table}

\subsection{A bateria de placebos}

O que sustenta a leitura da Seção~\ref{ap:shiftshare} não é a magnitude do
coeficiente, e sim a \emph{especificidade}: cada teste que deveria sair vazio
sai vazio. A Tabela~\ref{tab:placebos} traz a bateria inteira. Os placebos de
\emph{desfecho} trocam a variável dependente por outra que o câmbio não deveria
mover pela via testada --- área urbana e água; os de \emph{tempo} adiantam o
\emph{shifter} em um e dois anos, de modo que o efeito precederia a causa.

Convém dizer o que a bateria \emph{não} faz, porque é a ressalva que a
Seção~\ref{ap:shiftshare} desenvolve: todos estes são placebos de desfecho e de
tempo, e nenhum pergunta se a \emph{exposição} escolhida é a certa. Essa
pergunta é a da corrida entre exposições, e a resposta dela reduz o alcance da
frente.


\begin{table}[htb]
\centering
\caption[Placebos de desfecho e de tempo]{Bateria de placebos de desfecho e de tempo do desenho \emph{shift-share}.}
\label{tab:placebos}
\small
\begin{tabular}{lp{4.2cm}lcccc}
\toprule
& \textbf{Placebo} & \textbf{Tipo} & \textbf{$\hat{\gamma}$} & \textbf{EP} & \textbf{$p$} & \textbf{$R^2_w$} \\
\midrule
\multicolumn{7}{l}{\textit{H1 proxy de área (\#38)}} \\
 & $\Delta$ Área urbana (placebo) & desfecho & $+0{,}0029$ & $0{,}0045$ & $0{,}520$ & $0{,}00003$ \\
 & $\Delta$ Água (placebo) & desfecho & $-0{,}0342$ & $0{,}0294$ & $0{,}244$ & $0{,}00123$ \\
\multicolumn{7}{l}{\textit{H2 aptidão exógena (\#52)}} \\
 & $\Delta$ Área urbana (placebo) & desfecho & $-0{,}0040$ & $0{,}0042$ & $0{,}343$ & $0{,}00006$ \\
 & $\Delta$ Água (placebo) & desfecho & $+0{,}0272$ & $0{,}0281$ & $0{,}332$ & $0{,}00078$ \\
\multicolumn{7}{l}{\textit{H1 proxy de área (\#38)}} \\
 & nan & tempo & $-0{,}0177$ & $0{,}0095$ & $0{,}062$ & $0{,}00033$ \\
 & nan & tempo & $-0{,}0246$ & $0{,}0136$ & $0{,}070$ & $0{,}00065$ \\
\multicolumn{7}{l}{\textit{H2 aptidão exógena (\#52)}} \\
 & nan & tempo & $+0{,}0243$ & $0{,}0150$ & $0{,}105$ & $0{,}00062$ \\
 & nan & tempo & $+0{,}0045$ & $0{,}0154$ & $0{,}772$ & $0{,}00002$ \\
\bottomrule
\end{tabular}
\fontefig{elaboração própria.}
\notafig{mesma equação da Seção~\ref{ap:shiftshare}, com o desfecho ou a defasagem trocados. $n = 6.308$ observações, 166 unidades; efeitos fixos de unidade e de ano; erros-padrão agrupados nas duas dimensões. Um placebo cumpre o seu papel quando \textbf{não} rejeita: aqui nenhum dos oito o faz.}
\end{table}

\subsection{A grade confirmatória completa}

A Tabela~\ref{tab:confirmatorio} traz \emph{todas} as combinações de exposição,
desfecho e defasagem que a grade confirmatória rodou, e não apenas as
comentadas no Capítulo~\ref{cap:resultados}. A distinção entre grade
confirmatória e exploratória foi fixada antes da estimação, e o controle de
multiplicidade descrito na Seção~\ref{sec:met-instrumental} incide sobre a
segunda. Os desfechos de área --- pastagem e agricultura --- são nulos, e isso
consta aqui com o mesmo destaque dos demais.


\begin{table}[htb]
\centering
\caption[Grade confirmatória do \emph{drive} comum]{Grade confirmatória completa do desenho \emph{shift-share}.}
\label{tab:confirmatorio}
\small
\begin{tabular}{p{3.2cm}p{1.9cm}cccccc}
\toprule
\textbf{Exposição} & \textbf{Desfecho} & \textbf{Def.} & \textbf{Sinal} & \textbf{$\hat{\gamma}$} & \textbf{EP} & \textbf{$p$} & \textbf{$R^2_w$} \\
\midrule
Fronteira (veg. convertível, Norte) & $\Delta$ Rebanho & 0 & + & $+0{,}0147$ & $0{,}0138$ & $0{,}288$ & $0{,}00023$ \\
Fronteira (veg. convertível, Norte) & $\Delta$ Rebanho & 1 & + & $+0{,}0285$ & $0{,}0132$ & $0{,}031$ & $0{,}00089$ \\
Fronteira (veg. convertível, Norte) & $\Delta$ Pastagem & 0 & + & $-0{,}0111$ & $0{,}0355$ & $0{,}755$ & $0{,}00017$ \\
Fronteira (veg. convertível, Norte) & $\Delta$ Pastagem & 1 & + & $+0{,}0067$ & $0{,}0328$ & $0{,}838$ & $0{,}00007$ \\
Aptidão agrícola (Sul) & $\Delta$ Agricultura & 0 & + & $-0{,}0319$ & $0{,}0471$ & $0{,}498$ & $0{,}00175$ \\
Aptidão agrícola (Sul) & $\Delta$ Agricultura & 1 & + & $-0{,}0305$ & $0{,}0486$ & $0{,}530$ & $0{,}00149$ \\
Especialização em pasto & $\Delta$ Pastagem & 0 & - & $-0{,}0036$ & $0{,}0302$ & $0{,}906$ & $0{,}00002$ \\
Especialização em pasto & $\Delta$ Pastagem & 1 & - & $+0{,}0059$ & $0{,}0284$ & $0{,}835$ & $0{,}00005$ \\
Aptidão edafoclim. (exógena) & $\Delta$ Rebanho & 0 & - & $-0{,}0016$ & $0{,}0118$ & $0{,}890$ & $0{,}00000$ \\
Aptidão edafoclim. (exógena) & $\Delta$ Rebanho & 1 & - & $-0{,}0325$ & $0{,}0146$ & $0{,}026$ & $0{,}00116$ \\
Aptidão edafoclim. (exógena) & $\Delta$ Agricultura & 0 & + & $-0{,}0062$ & $0{,}0210$ & $0{,}770$ & $0{,}00006$ \\
Aptidão edafoclim. (exógena) & $\Delta$ Agricultura & 1 & + & $-0{,}0111$ & $0{,}0213$ & $0{,}603$ & $0{,}00020$ \\
Aptidão edafoclim. (exógena) & $\Delta$ Pastagem & 0 & - & $+0{,}0018$ & $0{,}0175$ & $0{,}920$ & $0{,}00000$ \\
Aptidão edafoclim. (exógena) & $\Delta$ Pastagem & 1 & - & $-0{,}0026$ & $0{,}0123$ & $0{,}832$ & $0{,}00001$ \\
\bottomrule
\end{tabular}
\fontefig{elaboração própria.}
\notafig{\emph{shifter}: variação padronizada do índice de taxa de câmbio efetiva real. ``Def.'' é a defasagem em anos do \emph{shifter}; ``Sinal'' é a direção prevista antes da estimação. Efeitos fixos de unidade e de ano; erros-padrão agrupados nas duas dimensões; 166 unidades. O $p$ desta tabela é o do erro-padrão agrupado, que a Seção~\ref{sec:met-instrumental} mostra ser otimista neste desenho. A régua de permutação não foi computada para estas linhas: ela existe para a especificação confirmatória do rebanho, na Tabela~\ref{tab:horserace}, cujas regressões são outras. Os $p$ abaixo, portanto, são limite inferior da incerteza, e nenhum deles é lido como significância no corpo do texto.}
\end{table}

\section{Teto de oferta: disponibilidade e depleção}
\label{ap:oferta}

A frente do teto de oferta separa duas afirmações que o mesmo painel testa. A
primeira é de \emph{nível}: o fluxo de conversão escala com o estoque ainda
disponível. A segunda é de \emph{comportamento}: a taxa de conversão cai à
medida que a unidade se deplete --- e é essa que distingue esgotamento de mera
aritmética, porque um fluxo proporcional ao estoque cairia sozinho sem que a
taxa mudasse.

A segunda afirmação é a que a decisão D29 reabriu. A variável de depleção,
construída como \(1 - \text{estoque}_t/\text{estoque}_{1985}\), sai do intervalo
\([0,1]\) em cerca de 14\% do painel, nas unidades que \emph{ganharam} estoque.
Padronizar uma variável com domínio violado dilui o sinal, e o nulo publicado
era artefato disso. A Tabela~\ref{tab:deplecao} traz a grade inteira de
tratamentos dessa variável, com e sem ponderação pelo tamanho do estoque, em
duas amostras --- e não apenas a combinação adotada.


\begin{table}[htb]
\centering
\caption[Canal de disponibilidade]{O canal de disponibilidade: fluxo de conversão contra estoque.}
\label{tab:disponibilidade}
\footnotesize
\begin{tabular}{p{4.0cm}p{2.6cm}ccccr}
\toprule
\textbf{Especificação} & \textbf{Regressor} & \textbf{$\hat{\beta}$} & \textbf{EP} & \textbf{$p$} & \textbf{$R^2_w$} & \textbf{$n$} \\
\midrule
B1  fluxo $\sim$ estoque\_def (disponibilidade) & log estoque (z) & $+2{,}756$ & $0{,}472$ & $<0{,}001$ & $0{,}434$ & 6.396 \\
B1q fluxo $\sim$ estoque + estoque$^2$ (não-linear) & log estoque (z) & $+2{,}686$ & $0{,}476$ & $<0{,}001$ & $0{,}434$ & 6.396 \\
B1q fluxo $\sim$ estoque + estoque$^2$ (não-linear) & log estoque$^2$ (z) & $+0{,}045$ & $0{,}451$ & $0{,}920$ & $0{,}434$ & 6.396 \\
B2a hazard $\sim$ estoque\_def & log estoque (z) & $-0{,}319$ & $0{,}189$ & $0{,}092$ & $-0{,}009$ & 6.379 \\
B2b hazard $\sim$ depleção\_def & log depleção (z) & $-0{,}015$ & $0{,}022$ & $0{,}481$ & $-0{,}000$ & 6.379 \\
B3  fluxo $\sim$ estoque + demanda (sem ano FE) & log estoque (z) & $+2{,}763$ & $0{,}391$ & $<0{,}001$ & $0{,}436$ & 6.396 \\
B3  fluxo $\sim$ estoque + demanda (sem ano FE) & $\Delta$ câmbio real (z) & $+0{,}022$ & $0{,}007$ & $0{,}003$ & $0{,}436$ & 6.396 \\
B3  fluxo $\sim$ estoque + demanda (sem ano FE) & $\Delta$ preço da soja (z) & $-0{,}030$ & $0{,}006$ & $<0{,}001$ & $0{,}436$ & 6.396 \\
B3  fluxo $\sim$ estoque + demanda (sem ano FE) & $\Delta$ crédito rural (z) & $+0{,}001$ & $0{,}005$ & $0{,}828$ & $0{,}436$ & 6.396 \\
\bottomrule
\end{tabular}
\fontefig{elaboração própria.}
\notafig{desfecho: fluxo anual de conversão. Efeitos fixos de unidade e de ano; erros-padrão agrupados nas duas dimensões. Regressores em \emph{escore-z}.}
\end{table}

\begin{table}[htb]
\centering
\caption[Grade de tratamentos da variável de depleção]{Grade completa de tratamentos da variável de depleção (decisão D29).}
\label{tab:deplecao}
\footnotesize
\begin{tabular}{p{4.0cm}llcccc}
\toprule
\textbf{Tratamento da depleção} & \textbf{Amostra} & \textbf{Peso} & \textbf{$\hat{\beta}$} & \textbf{$p_{ent}$} & \textbf{$p_{ent+ano}$} & \textbf{$R^2_w$} \\
\midrule
sem tratamento (a do \#39) & todas & --- & $-0{,}015$ & $0{,}317$ & $0{,}481$ & $-0{,}0002$ \\
sem tratamento (a do \#39) & todas & estoque & $-1{,}165$ & $0{,}050$ & $0{,}060$ & $0{,}1043$ \\
piso em 0 (não descarta linha) & todas & --- & $-0{,}202$ & $<0{,}001$ & $0{,}053$ & $0{,}0451$ \\
piso em 0 (não descarta linha) & todas & estoque & $-0{,}215$ & $<0{,}001$ & $<0{,}001$ & $0{,}1846$ \\
domínio documentado [0,1] & todas & --- & $-0{,}311$ & $<0{,}001$ & $0{,}039$ & $0{,}0921$ \\
domínio documentado [0,1] & todas & estoque & $-0{,}202$ & $<0{,}001$ & $<0{,}001$ & $0{,}1947$ \\
winsor p1 (não põe no domínio) & todas & --- & $-0{,}069$ & $0{,}026$ & $0{,}111$ & $0{,}0014$ \\
winsor p1 (não põe no domínio) & todas & estoque & $-0{,}497$ & $<0{,}001$ & $<0{,}001$ & $0{,}1661$ \\
sem tratamento (a do \#39) & corte1k & --- & $-0{,}291$ & $0{,}242$ & $0{,}262$ & $0{,}0298$ \\
sem tratamento (a do \#39) & corte1k & estoque & $-1{,}097$ & $<0{,}001$ & $<0{,}001$ & $0{,}1819$ \\
piso em 0 (não descarta linha) & corte1k & --- & $-0{,}440$ & $<0{,}001$ & $<0{,}001$ & $0{,}1341$ \\
piso em 0 (não descarta linha) & corte1k & estoque & $-0{,}400$ & $<0{,}001$ & $<0{,}001$ & $0{,}1879$ \\
domínio documentado [0,1] & corte1k & --- & $-0{,}452$ & $<0{,}001$ & $<0{,}001$ & $0{,}1722$ \\
domínio documentado [0,1] & corte1k & estoque & $-0{,}397$ & $<0{,}001$ & $<0{,}001$ & $0{,}1968$ \\
winsor p1 (não põe no domínio) & corte1k & --- & $-0{,}444$ & $0{,}002$ & $0{,}009$ & $0{,}0664$ \\
winsor p1 (não põe no domínio) & corte1k & estoque & $-0{,}808$ & $<0{,}001$ & $<0{,}001$ & $0{,}1854$ \\
nan & todas & --- & $+2{,}756$ & $<0{,}001$ & $<0{,}001$ & $0{,}4341$ \\
nan & todas & --- & $+2{,}686$ & $<0{,}001$ & $<0{,}001$ & $0{,}4342$ \\
nan & todas & --- & $-0{,}320$ & $0{,}002$ & $0{,}092$ & $-0{,}0095$ \\
nan & corte1k & --- & $+2{,}690$ & $<0{,}001$ & $<0{,}001$ & $0{,}4342$ \\
nan & corte1k & --- & $+2{,}623$ & $<0{,}001$ & $<0{,}001$ & $0{,}4342$ \\
nan & corte1k & --- & $-0{,}057$ & $0{,}632$ & $0{,}737$ & $-0{,}0034$ \\
\bottomrule
\end{tabular}
\fontefig{elaboração própria.}
\notafig{desfecho: taxa anual de conversão. $p_{ent}$ agrupa por unidade e $p_{ent+ano}$ nas duas dimensões, que é a régua adotada no teste de origem e a reportada no corpo do texto. ``Peso'' indica ponderação pelo tamanho do estoque. A linha ``sem tratamento'' é a especificação publicada antes da D29, e o nulo dela é o artefato que a decisão corrigiu.}
\end{table}

\section{Desenvolvimento municipal e o gradiente}
\label{ap:ifdm}

A última conta do Capítulo~\ref{cap:resultados} pergunta se a expansão da área
cultivada comprou desenvolvimento. Ela tem dois blocos, e os dois são
\textbf{associativos}: nenhum identifica efeito causal.

O primeiro é transversal, sobre os 246 municípios, e mede se o índice de
desenvolvimento se ordena pela latitude --- em nível e em variação --- com e sem
controle de longitude, que separa o gradiente sul-norte do gradiente
leste-oeste. O segundo é um painel de efeitos fixos \emph{dentro} do município,
que pergunta quanto de desenvolvimento acompanha uma duplicação da área ou do
valor adicionado.


\begin{table}[htb]
\centering
\caption[Gradiente do desenvolvimento municipal]{Gradiente latitudinal do desenvolvimento municipal, com e sem controle de longitude.}
\label{tab:ifdm-gradiente}
\footnotesize
\begin{tabular}{llcccccc}
\toprule
\textbf{Desfecho} & \textbf{Spec.} & \textbf{$n$} & \textbf{$R^2$} & \textbf{$\hat{\beta}_{lat}$} & \textbf{$p$} & \textbf{$\hat{\beta}_{lon}$} & \textbf{$p$} \\
\midrule
IFDM 2023 (nível) & $\sim$lat & 246 & $0{,}246$ & $-0{,}0383$ & $<0{,}001$ & --- & --- \\
IFDM 2023 (nível) & $\sim$lat+lon & 246 & $0{,}252$ & $-0{,}0361$ & $<0{,}001$ & $-0{,}0059$ & $0{,}213$ \\
$\Delta$ IFDM 2013$\rightarrow$2023 & $\sim$lat & 245 & $0{,}022$ & $+0{,}0078$ & $0{,}018$ & --- & --- \\
$\Delta$ IFDM 2013$\rightarrow$2023 & $\sim$lat+lon & 245 & $0{,}025$ & $+0{,}0066$ & $0{,}066$ & $+0{,}0032$ & $0{,}415$ \\
\bottomrule
\end{tabular}
\fontefig{elaboração própria.}
\notafig{mínimos quadrados sobre a seção transversal dos municípios; latitude e longitude em graus decimais do centroide. Um $\hat{\beta}_{lat}$ negativo no nível significa índice menor mais ao norte. A leitura é associativa.}
\end{table}

\begin{table}[htb]
\centering
\caption[Painel de efeitos fixos do desenvolvimento]{Desenvolvimento municipal em painel de efeitos fixos, dentro do município.}
\label{tab:ifdm-painel}
\footnotesize
\begin{tabular}{llccc}
\toprule
\textbf{Modelo} & \textbf{Regressor} & \textbf{$\hat{\beta}$} & \textbf{$p$} & \textbf{$R^2_w$} \\
\midrule
d\_l\_va & $\Delta$ log VA agro & $+0{,}0039$ & $0{,}539$ & $0{,}00049$ \\
d\_l\_area & $\Delta$ log área cultivada & $+0{,}0118$ & $0{,}011$ & $-0{,}00033$ \\
d\_l\_va+d\_l\_area & $\Delta$ log VA agro & $+0{,}0035$ & $0{,}582$ & $0{,}00034$ \\
d\_l\_va+d\_l\_area & $\Delta$ log área cultivada & $+0{,}0117$ & $0{,}012$ & $0{,}00034$ \\
\bottomrule
\end{tabular}
\fontefig{elaboração própria.}
\notafig{variáveis em log-diferença; efeitos fixos de município e de ano. \texttt{d\_l\_area} é a variação log da área cultivada e \texttt{d\_l\_va} a do valor adicionado agropecuário. O $R^2$ \emph{within} negativo indica ajuste pior que a média dentro do município, e é reportado como saiu.}
\end{table}
